\scanpage{44}
% archive sheet number: 43
Revenons aux conditions de la \uline{proposition}, et \unsure{déterminons|déterminer}
\[
\omega_{G_0}=t_{G_0}^{\vee}\simeq M_{\unsure{\nu|v}}/\pi^{\unsure{F|f|r}}M_\nu\simeq \Omega/p\Omega
\]
avec $\deleted{\unsure{K}}k=A/\pi A$ y opérant de la façon évidente, et la structure de vectoriel sur $\inserted{W/pW=}k=\mathbb F_q$ est déduite de la structure précédente grâce à $\sigma^\nu$. Le cas le plus simple est celui où $\nu=0$, i.e. où \deleted{\unsure{les deux}} les deux opérations de $k$ sur $\omega_{G_0}$ coïncident. \deleted{\unsure{Si on sup|Soit}} (les autres cas s'en déduisent grâce à \unsure{autom} $\sigma^\nu$ de $k=\mathbb F_q$ \deleted{\unsure{sur $k_p$}}).

La filtration canonique sur $M\otimes_W k=\sum M_\nu\otimes_A k$ donne comme sous-module (de rang $1$) \deleted{\unsure{$M_{f-1}\otimes_A k$}} \inserted{$\omega_0$ \deleted{\unsure{$=$}} tel que $\omega_0^{(p)}\subset(M\otimes_W k)^{(p)}\unsure{=}$ soit $M_{f-1}\otimes_A k$ ce qui donne $\omega_0=M_0\otimes_A\deleted{\unsure{k}}k$.}

Le cas \unsure{standard|standart} est
\[
G_0=G_0(A,\pi)\quad\text{correspondant à }\;(\nu=0\text{ et })\;\Omega=A
\]
on en déduit le cas d'un $\Omega$ quelconque par
\[
\deleted{\unsure{G_0}}\quad\Omega\otimes_A G_0(A)
\]

Donc on a dans ce cas comme module de Dieudonné
\[
\left\{\begin{aligned}
M&=A\otimes_{\mathbb Z_p}W\\
F_M&=\mathrm{id}_A\otimes_{\mathbb Z_p}\sigma+\operatorname{pr}_0\bigl((\pi-1)\otimes\sigma\bigr)
\end{aligned}\right.
\]
où $\operatorname{pr}_0$ est la projection sur le composant correspondant au $A$-hom $A\otimes_{\mathbb Z_p}W\longrightarrow A$ qui sur $W$ est l'inclusion $W\mathrel{\overset{i}{\unsure{\lhook\joinrel\relbar|\hookrightarrow}}}A$.

\begin{center}\rule{2.2em}{0.3pt}\\[-4pt]\rule{1.2em}{0.3pt}\end{center}

Nous nous proposons maintenant de relever $G_0$ \deleted{l'action} en un groupe de BT $G$ sur $A$, \inserted{comm.} avec action de $A$ sur $G$, de telle façon que l'action de $A$ sur $t_G$ (qui est un $A$-module inversible) soit l'action évidente. D'après Lubin-Tate,
